\chapter{Single-Raft and ReplayPool}

Multi-Raft is the direct replacement for Paxos, but Single-Raft is the more
interesting systems design. It reduces replication structure by placing all
partitions in one Raft group per process. That consolidation is attractive,
but it creates a new pressure point: one ordered log can feed one apply path
faster than follower database replay can consume entries.

The thesis separates three concepts that are easy to conflate:

\begin{table}[htbp]
\centering
\caption[Single-Raft pipeline stages]{Single-Raft pipeline stages and their roles.}
\label{tab:replay-stages}
\small
\begin{tabularx}{\linewidth}{@{}XX@{}}
\toprule
Stage & Meaning \\
\midrule
Raft commit & consensus decides that a log entry is safe \\
Apply staging & \code{raft\_apply} dispatches committed entries into Mako \\
Follower replay & Mako replays transaction effects into database state \\
\bottomrule
\end{tabularx}
\end{table}

The final implementation does not make \code{raft\_apply} perform all heavy replay.
Instead, \code{raft\_apply} is kept lightweight. It stages committed entries and
invokes callbacks. The Mako callback copies the committed payload and enqueues
heavy work into ReplayPool. Replay workers then execute follower-side replay
in parallel.

\section{The Bottleneck}

Without replay offload, Single-Raft risks a fan-in bottleneck. Many Mako
workers can submit work, and one Raft group can order the work, but the
apply/replay path can become narrow. As worker count increases, follower
replay work accumulates. A benchmark may report high leader-side throughput
while followers make much less replay progress.

\begin{figure}[htbp]
\centering
\includegraphics[width=0.96\linewidth]{figures/manual/fig04_single_raft_replay_bottleneck.pdf}
\caption[Single-Raft replay bottleneck]{Single-Raft replay bottleneck before ReplayPool. Many worker threads feed one ordered Raft log, but follower replay can accumulate backlog without replay offload.}
\label{fig:single-raft-replay-bottleneck}
\end{figure}

The important claim is not that Single-Raft is fundamentally slow. The claim
is that Single-Raft changes where parallelism must be preserved. If replay is
kept serial, Single-Raft can look healthy from leader throughput alone while
follower state lags behind. If replay is parallelized, the design can recover
useful replicated progress.

\section{ReplayPool Design}

ReplayPool moves heavy follower-side replay out of the Raft apply staging
thread. The pipeline becomes Raft commit \(\rightarrow\) \code{raft\_apply}
staging \(\rightarrow\) Mako callback \(\rightarrow\) ReplayPool queue
\(\rightarrow\) replay workers.

The benchmark configuration called 1:1 means that
\code{MAKO\_REPLAY\_THREADS} equals the Mako worker-thread count.

It does not mean one Raft group per worker. It does not mean one process per
worker. It does not mean one apply thread per worker.

\begin{figure}[htbp]
\centering
\includegraphics[width=0.96\linewidth]{figures/manual/fig05_replay_pool_pipeline.pdf}
\caption[ReplayPool pipeline]{Final Single-Raft replay pipeline. The Raft apply path stages committed payloads, while ReplayPool workers perform heavy follower-side replay in parallel.}
\label{fig:replaypool-pipeline}
\end{figure}

\section{Why ReplayPool Is a Systems Contribution}

Consensus protocols are often evaluated as if state-machine application cost
were small. Mako makes that assumption unsafe. The replicated command is a
transaction log, and follower replay updates database state. That work can be
substantial.

ReplayPool is therefore not a cosmetic optimization. It is the mechanism that
allows Single-Raft to preserve Mako's multicore structure after consolidating
the consensus log. The evaluation supports this by comparing leader-side
throughput with follower replay progress and by showing that larger replay
pools let follower replay catch up to the leader-side rate.
